\begin{claim}{AG.34}{Proved; nontrivial primitive-cover counterexamples over $\mathbf F_{2^{38}}$}
For the family $C_t$ of AG.33 over $\mathbf F_{2^{38}}$,
at least $59927733332$ parameters $t$ define curves having
no nontrivial connected cyclic \'etale cover of odd degree
with primitive two-rank zero. This counts parameters,
without asserting that the corresponding curves are
pairwise nonisomorphic or selecting one individual parameter.
The degree-one qualification at the end of AG.33 applies.
\end{claim}
\begin{proof}
There is a useful refinement of the torsion incidence bound.
For an integer $e\ge1$ let
\[
\begin{split}
 \mathcal S_e={}&\{n:\ 3\le n<10e^3,\ n\text{ odd},\
       M_e(n)(M_e(n)+1)\le2(n-1)\},\\
 M_e(n)={}&\frac{\varphi(n)}{\gcd(\operatorname{ord}_n(2),e)},\qquad
 J_4(n)=n^4\prod_{\ell\mid n}(1-\ell^{-4}).
\end{split}                                                   \tag{AG.34.1}
\]
The integer $J_4(n)$ is the number of points of exact order
$n$ in $A[n]\simeq(\mathbf Z/n\mathbf Z)^4$.
For $n\in\mathcal S_e$ define $\mathcal B_{e,n}$ to be
the parameters $t\in\mathbf F_{4^e}\setminus\mathbf F_4$
for which the pencil member $D_t$ contains at least
$M_e(n)$ distinct points of exact order $n$ in $A$.
Every nonzero point belongs to at most one pencil member.
Counting incidences therefore gives
\[
               \#\mathcal B_{e,n}
                  \le\left\lfloor\frac{J_4(n)}{M_e(n)}\right\rfloor.
                                                               \tag{AG.34.2}
\]
If $C_t$ admits a nontrivial primitive-zero cover of
conductor $n$, AG.32 puts $n$ in $\mathcal S_e$ and gives
at least $M_e(n)$ exact-order-$n$ points of its theta
divisor. The odd-torsion isomorphism induced by $v_t$
sends them injectively to points counted in
$\mathcal B_{e,n}$. Thus the explicitly defined set
\[
 \mathbf F_{4^e}\setminus
       \left(\mathbf F_4\cup
          \bigcup_{n\in\mathcal S_e}\mathcal B_{e,n}\right)
                                                               \tag{AG.34.3}
\]
consists entirely of parameters with the asserted property,
and its size is at least
\[
       4^e-4-
       \sum_{n\in\mathcal S_e}
          \left\lfloor\frac{J_4(n)}{M_e(n)}\right\rfloor.
                                                               \tag{AG.34.4}
\]
The definition of $\mathcal B_{e,n}$ uses the threshold
$M_e(n)$; the stronger bound (AG.34.2) is not being asserted
for parameters with only one theta point of order $n$.

For $e=19$, exact enumeration below the proved global
cutoff $68590$ gives precisely
\[
 \mathcal S_{19}=\{191,229,361,419,457,571,573,647,687,
                         955,1083,1257,1371\}.
\]
The certificate records all factorizations, totients,
multiplicative orders, values of $M_{19}(n)$, exact-order
point counts, and integer quotients. Its sum is
\[
 \sum_{n\in\mathcal S_{19}}
      \left\lfloor\frac{J_4(n)}{M_{19}(n)}\right\rfloor
       =214950173608.
\]
As $4^{19}=274877906944$, formula (AG.34.4) becomes
\[
             274877906944-4-214950173608
                     =59927733332>0.
\]
The root program \path{code/A6_scalar_conductor_bounds.py}
and its complete JSON output give the global enumeration.
The independent review
\path{agent_geometric/code/primitive_counterexample_e19_review.py}
factors every odd integer below the cutoff by trial division
and recomputes every relevant multiplicative order without
using the root program's sieve or cache. Its complete JSON
agrees on every $e=19$ record. These are exhaustive finite
arithmetic certificates for the displayed existence theorem,
not an enumeration of a parameter in (AG.34.3).
\end{proof}
